\section{The Stupid Twenty First Century}

Rather than watch the Academy Awards, I had some time to glance through some ``conservative'' or ``right'' web sites. The striking thing I noticed is the trend toward an Oprah Winfrey type of ``confessional'' writing and ``conversion'' experiences. On this, we agree with H L Mencken who never trusted anyone who had too many conversions. My recommendation is to end such writings immediately: it is demeaning and humiliating.

The other thing to notice is the complete lack of agreement. The Left is united on goals, just not, perhaps, on the method or rapidity of change. The alleged right is all over the place. Everyone loves a winner, and the Left is winning. It seems unstoppable, on the winning side of history, and without credible opposition.

\textbf{Leon Daudet} was a co-founder, along with Charles Maurras, of the journal \emph{Action Française} and sympathetic to Rene Guenon (who likewise was, for a time, sympathetic to the goals of Action Francaise.) Guenon even refers to Daudet more than once. Regarding Guenon, Daudet wrote:

\begin{quotex}
Since the time of the Encyclopedists, or even the Reformation, the West has been in a state of intellectual anarchy that is a veritable barbarism.

The civilization of which it is so proud rests upon a collection of material and industrial improvements (which exacerbate the risk of war) and a weak moral and intellectual substructure which lacks any metaphysical basis.

By different paths, I had reached a similar conclusion in my own analysis of the benighted nineteenth century, but my ignorance of Eastern philosophy, of which Guenon is a master, prevented me from laying out the formidable parallel that he presents to us. Although he does not expressly say so, it is evident that the West is threatened more from within, on account of its mental weakness, than from outside, where its prospects nevertheless are also not so certain.
\end{quotex}


So those who suddenly have become aware of a `problem' in the West, are really johnny-come-latelys. Perhaps the United States were immune from the effects of the European malaise and are just now catching up. Daudet traces it back to the 19th century, but Guenon and Evola, detect its roots in the Reformation, Enlightenment, and Revolution. Those who see in those historical movements the solution to the problem are unfortunately mistaken, despite their sincerity and good intentions. They have no place on the ``right''. In his book, published in 1922, \emph{The Stupid 19th Century}, Daudet writes:


\begin{quotex}
The liberal is a man who venerates the good God but respects the devil He aspires to order and praises anarchy. ... The liberal dominates the nineteenth century. He is its teacher and pride. .. on the social level, as on the economic and political level, the harmful effects of liberalism are innumerable.

The unjustly scoffed at, vilified, execrated term ``Reaction'' must be boldly taken up again, if we want to chase away the bloodthirsty error, if we lead back to the true peace, institutions and notions of life with the collapse of the notions and institutions of death, honored in the 19th century.
\end{quotex}


Clearly, the collapse of the ideas of the 19th century has not yet occurred. Should it ever happen, it will be from its own inner contradictions, not from any intellectual threat from the right. Daudet leaves us with a list of deadly ideas from that century, though they are just as relevant today. Use it as a self-test to determine which side you are really on. I think a score of 20 or higher should suffice.

\paragraph{The clichés or deadly ideas of the 19th century}

\begin{enumerate}


\item The 19th century is the century of science.

\item The 19th century is the century of progress.

\item The 19th century is the century of democracy, which is progress and progress continues.

\item The darkness of the Middle Ages.

\item The Revolution is sacred and emancipated the French people.

\item Democracy is peace. If you want peace, prepare the peace.

\item The future is science. Science is always beneficial\footnote{Global warming?} 

\item Secular education is the emancipation of the people\footnote{Check out the test scores.} 

\item Religion is the child of fear.

\item It is States that fight. The people are always ready to come to agreement.

\item It is necessary to replace the study of Latin and Greek, which have become useless, with those living languages which are useful.

\item The relations between peoples are ceaselessly going to improve. We are racing toward a United States of Europe (EU).

\item Science has neither borders nor country.

\item The people thirst for equality\footnote{Some are more equal than others.}.

\item We are at the dawn of a new era of brotherhood and justice.

\item Property is theft. Capital is war.

\item All religions are the same, as long as they admit to the divine\footnote{And universal brotherhood.}.

\item God exists only in and through human consciousness. This consciousness creates God a little more each day\footnote{I am god.}. 

\item Evolution is the law of the universe\footnote{Every day I am getting better.}. 

\item Men are born naturally good. It is society that corrupts them.

\item There are only relative truths; absolute truth does not exist\footnote{Except this one.}. 

\item All opinions are good and worthy.
\end{enumerate}



\flrightit{Posted on 2011-02-27 by Cologero}

\begin{center}* * *\end{center}

\begin{footnotesize}
\begin{sffamily}

\texttt{James O'Meara on 2011-02-28 at 19:39 said:}

In Mann's Dr. Faustus, Leverkuhn says something like: ``The nineteen century must have been an unusually comfortable place; no one seems to want to leave it.''

\end{sffamily}
\end{footnotesize}

